\section{Datasets}
\label{sec:datasets}
Web-scale corpora such as Koala-36M \cite{koala36m2024}, WebVid-10M \cite{webvid10m2021}, Panda-70M \cite{panda70m2024} and HD-VG-130M \cite{hdvg130m2024} collectively exceed 250 M clips, yet their stock-footage provenance yields noisy captions and licences that forbid commercial use. High-definition, human-centric datasets like CelebV-HQ \cite{celebvhq2022}, OpenHumanVid \cite{openhvid2024}, HumanVid \cite{humanvid2024} and HDTF \cite{hdtf2021} supply face tracks, skeletons and camera-motion labels, but most clips remain under 20 s, limiting long-form training.

Vript \cite{vript2024} provides six-minute films with 145-word scene-level “scripts.” Large-scale video–text datasets like HD-VILA-100M \cite{hd-vila2022} and Panda-70M \cite{panda70m2024} have enabled text-to-video generation at scale, but their clips are mostly 5–15 s with minimal narrative context. To push toward longer, story-driven videos, recent benchmarks offer richer structure: MiraData provides 1–2 min sequences with dense, structured captions covering objects, actions, style and camera motions \cite{miradata2024}, and MovieBench is the first movie-level dataset with hierarchical annotations (movie, scene, shot) enforcing character consistency and multi-scene storytelling \cite{moviebench2025}. Examples of these datasets used in video generation models are in Table~\ref{tab:tasks-datases}.